\documentclass{article}
\usepackage[T1]{fontenc}
\usepackage{lmodern}
\usepackage{graphicx}
\usepackage{amsmath,amssymb}
\usepackage[a4paper,margin=2cm]{geometry}
\usepackage[absolute,overlay]{textpos}
\setlength{\TPHorizModule}{1mm}
\setlength{\TPVertModule}{1mm}
\usepackage[indonesian]{babel}
\usepackage{hyperref}
\setlength{\emergencystretch}{2em}
% unit: Halaman 42

\begin{document}

% === Halaman 42 (llm) ===

Contoh: Caesar cipher (ingat kembali materi di dalam kuliah Matdis $\smiley$ )
\begin{itemize}
    \item $\mathcal{M} = \{ \text{huruf-huruf alfabet} \}$
    \item $\mathcal{K} = \{ k \mid k \text{ adalah bilangan bulat dan } 0 \le k \le 25 \}$
    \item $\mathcal{E} = \{ E_k \mid k \in \mathcal{K} \text{ dan untuk semua plainteks } m, \\
    [ E_k(m) = (m + k) \pmod{26} ]$
    \item $\mathcal{D} = \{ D_k \mid k \in \mathcal{K} \text{ dan untuk semua cipherteks } c, \\
    [ D_k(c) = (c - k) \pmod{26} ]$
    \item $\mathcal{C} = \mathcal{M}$
\end{itemize}

\end{document}
